Three pre-registered possibilities resolved against the more
optimistic hypothesis and are retained as findings:
\begin{itemize}
\item \textbf{Weather adds nothing in Austin}: $+0.001$ test MAE on top
of calendar features ($p = 0.995$). The signal that produces a 9.3\%
gain in New York produces none in Austin.
\item \textbf{Pooled training mostly hurts}: $+1.8\%$ (New York,
$p=0.038$) and $+6.6\%$ (Chicago, $p=0.001$) test MAE versus local
models; zero-shot transfer degrades MAE by $+44$--$147\%$ except in San
Francisco ($+3.6\%$). With six years of local history, cross-city
transfer is not beneficial in this benchmark.
\item \textbf{Quantile intervals are under-dispersed}: 90\% intervals
cover 76.8--81.2\%. The uncertainty-aware policy wins under scarcity
despite this, but the calibration deficit is real and unrepaired in
this paper.
\end{itemize}
We also record that the perfect-information myopic reference was
exceeded once (Austin, generous regime, proportional policy, 104.6\% of
gap closed), which is possible because the reference is per-day myopic
rather than horizon-optimal; this motivated the precise naming used in
the main text.
